% !TeX root = ../../Book3.tex
\chapter{自守形式与 \texorpdfstring{\(L\)}{L}-函数}
\label{b3:app:E}
\begin{chapterabstract}
模形式的 Fourier 系数同时满足解析约束与乘法关系。前者来自变换律和尖点衰减，后者来自 Hecke 算子。以全模群为可完整计算的模型，本附录建立 Hecke 系数公式、特征形式的 Euler 乘积以及 Mellin 函数方程；高层结构则用同余子群和 Fricke 对合说明。自守表示与 Langlands 纲领只介绍其对象及对应问题，不把尚未建立的一般理论当作已证工具。
\end{chapterabstract}
\begin{learninggoals}
\begin{itemize}
\item\label{b3:item:E-goal-1}区分权、层、特征与尖点条件。
\item\label{b3:item:E-goal-2}由 Hecke 算子的求和式计算 Fourier 系数。
\item\label{b3:item:E-goal-3}把特征值关系变为局部 Euler 因子。
\item\label{b3:item:E-goal-4}从虚轴 Mellin 积分推导解析延拓与函数方程。
\end{itemize}
\end{learninggoals}
前置为第二十三章的模形式与第二十一章的 Mellin 变换。附录 D 的 Dirichlet 特征用于描述高层的变换因子，但全模群的主要证明不依赖附录 D 的素数结论。

\section{层与尖点}
\label{b3:sec:E01}

对正整数 \(N\)，定义\term[tongyu-ziqun]{同余子群}[congruence subgroup]的三个基本例子
\[
 \Gamma_0(N)=\left\{\begin{pmatrix}a&b\\c&d\end{pmatrix}\in\mathrm{SL}_2(\mathbb Z):c\equiv0\pmod N\right\},
\]
\[
 \Gamma_1(N)=\{\gamma\in\Gamma_0(N):a\equiv d\equiv1\pmod N\},
\]
\[
 \Gamma(N)=\{\gamma\in\mathrm{SL}_2(\mathbb Z):\gamma\equiv I\pmod N\}.
\]
它们的指标有限，因为约化到有限矩阵群后的核包含 \(\Gamma(N)\)。设 \(k\in\mathbb Z\)、\(f:\mathfrak H\to\mathbb C\)，\(\gamma=\bigl(\begin{smallmatrix}a&b\\c&d\end{smallmatrix}\bigr)\in\mathrm{GL}_2(\mathbb R)\) 且 \(\det\gamma>0\)。对 \(z\in\mathfrak H\) 采用
\[
 (f|_k\gamma)(z)=(\det\gamma)^{k/2}(cz+d)^{-k}f(\gamma z).
\]
这个复合律可以在分母上逐项核对。记 \(j(\gamma,z)=cz+d\)，矩阵乘法给
\[
j(\alpha\beta,z)=j(\alpha,\beta z)j(\beta,z),\qquad
\det(\alpha\beta)=\det\alpha\det\beta.
\]
因权为整数，两个分母的 \(-k\) 次幂可以直接相乘；正行列式的实数幂也相乘。因此
\((f|_k\alpha)|_k\beta=f|_k(\alpha\beta)\)。在全模群和整数偶权 \(k\) 上，这就是第二十三章的变换律。
\ProseParagraphBreak
模 \(N\) 特征 \(\chi\) 的模形式满足 \(f|_k\gamma=\chi(d)f\)，\(\gamma\in\Gamma_0(N)\)，并在所有尖点全纯。矩阵 \(-I\) 给必要的奇偶条件 \(\chi(-1)=(-1)^k\)。尖点为 \(\mathbb Q\cup\{\infty\}\) 在子群下的轨道；若 \(\sigma\in\mathrm{SL}_2(\mathbb Z)\) 把无穷远送到所考察尖点，某个正整数 \(w\) 使 \(f|_k\sigma\) 在足够的平移周期下不变。它在局部参数 \(e^{2\pi\mathrm{i}z/w}\) 中无负次幂，就是该尖点的全纯性；常数项为零称在那里消失。带非平凡特征时可先增大周期使相位因子消失，再作此检验。
\ProseParagraphBreak
以 \(M_k(\Gamma_0(N),\chi)\)、\(S_k(\Gamma_0(N),\chi)\) 分别表示模形式与尖点形式空间。无穷远的展开不够保证其他尖点也全纯。例如 \(f(z)=j(z)\) 虽在上半平面全纯，却在无穷远有极点，因此不是权0模形式。这一例子在高层同样提醒我们：域内全纯与尖点条件是两项独立要求。

一个能区分两个尖点的例子更能说明检验的必要性。设 \(N>1\)，取
\[
R_N(z)=\frac{\Delta(Nz)}{\Delta(z)}.
\]
第二十三章已证明 \(\Delta\) 在上半平面无零点，因此 \(R_N\) 在域内全纯。若 \(\gamma\in\Gamma_0(N)\)，则
\(N\gamma z=\begin{psmallmatrix}a&Nb\\c/N&d\end{psmallmatrix}(Nz)\)，右边的矩阵属于全模群。分子、分母均乘以 \((cz+d)^{12}\)，故 \(R_N\) 是权0的不变函数。在无穷远，\(R_N=q^{N-1}(1+O(q))\)，没有极点；但在尖点0，取 \(\sigma=S\)，有
\[
(R_N|_0S)(z)=N^{-12}\frac{\Delta(z/N)}{\Delta(z)}
=N^{-12}q_0^{1-N}(1+O(q_0)),\qquad q_0=e^{2\pi\mathrm{i}z/N}.
\]
这里周期取 \(N\) 是因为 \(ST^NS^{-1}\) 的左下角为 \(-N\)，属于 \(\Gamma_0(N)\)。负指数 \(1-N\) 显示0处有极点。可见无穷远的良好展开确实不足以保证模形式的尖点条件。


\section{全模群上的 Hecke 算子}
\label{b3:sec:E02}

先取偶数 \(k\ge2\) 与全模群 \(\Gamma=\mathrm{SL}_2(\mathbb Z)\)。我们希望把不同伸缩和平移下的 \(f\) 组合起来，仍得到满足同一模变换律的函数。最简单的素数情形是
\[
 p^{k-1}f(pz)+\frac1p\sum_{b=0}^{p-1}f\left(\frac{z+b}p\right).
\]
第一项沿 \(z\mapsto pz\) 伸缩，后一组沿 \(z\mapsto(z+b)/p\) 伸缩并求遍平移代表。权因子补偿伸缩，求遍代表使换基只置换求和项。下面的整数矩阵证明将核验这些选择，而不把模性视为公式形式上自动成立。

设 \(k\ge2\) 为偶整数，\(f\in M_k(\mathrm{SL}_2(\mathbb Z))\)。对正整数 \(n\)，将 \(n\) 的不同分解 \(ad=n\) 一并纳入，定义\term[hecke-suanzi]{Hecke 算子}[Hecke operator]
\begin{equation}\label{b3:eq:E-hecke-definition}
 (T_nf)(z)=n^{k-1}\sum_{ad=n\atop a,d>0}d^{-k}
                      \sum_{b=0}^{d-1}f\left(\frac{az+b}{d}\right).
\end{equation}
\(T_1\) 是恒等算子。这里的正规化使第一 Fourier 系数直接等于原函数的第 \(n\) 个系数。

\begin{proposition}[保持模性与系数公式]\label{b3:prop:E-hecke-coefficients}
设 \(k\ge2\) 为偶整数，\(\Gamma=\mathrm{SL}_2(\mathbb Z)\)，\(T_n\) 为按\eqref{b3:eq:E-hecke-definition}定义的权 \(k\) Hecke 算子，其前因子为 \(n^{k-1}\)。对每个正整数 \(n\)，\(T_n\) 保持 \(M_k(\Gamma)\) 与 \(S_k(\Gamma)\)。若 \(f\in M_k(\Gamma)\)、\(q=e^{2\pi\mathrm iz}\)，且 \(f=\sum_{m\ge0}a_mq^m\)，则 \(T_nf=\sum b_mq^m\)，其中
\[
 b_m=\sum_{d\mid(m,n)}d^{k-1}a_{mn/d^2}\quad(m\ge1),
 \qquad b_0=\sigma_{k-1}(n)a_0.
\]
\end{proposition}
\begin{proof}
行列式为 \(n\) 的整矩阵按左乘 \(\Gamma\) 分类。用整数行变换把第一列化为 \((a,0)^t\)，使 \(a>0\)，再把右上角模 \(d=n/a\) 化到 \(0\le b<d\)，得到唯一代表 \(\begin{psmallmatrix}a&b\\0&d\end{psmallmatrix}\)。因此式~\eqref{b3:eq:E-hecke-definition}等于 \(n^{k/2-1}\) 乘所有左陪集代表的 \(f|_k\alpha\) 之和。右乘任一 \(\gamma\in\Gamma\) 只置换这些陪集，故和满足模变换律；各项在上半平面全纯。
将 Fourier 级数代入，对 \(b\) 求和时，\(\sum_{b=0}^{d-1}e^{2\pi\mathrm{i}rb/d}\) 仅在 \(d\mid r\) 时等于 \(d\)，其余为零。剩余项为
\[
 \sum_{ad=n}a^{k-1}\sum_{\ell\ge0}a_{d\ell}q^{a\ell}.
\]
提取 \(q^m\)，即有 \(a\mid m\) 且系数为 \(a^{k-1}a_{nm/a^2}\)，得到公式。没有负次幂，故尖点全纯；若 \(a_0=0\)，常数项也为零。全模群仅有一个尖点轨道，因此无需再检验其他轨道。
\end{proof}
矩阵分类在几何上是有限指标子格的分类；Hecke 算子把一条椭圆曲线上所有相应次数的同源关系合并求和。这里的证明只使用整数矩阵，不需要先建立一般同源理论。

\begin{proposition}[乘法关系]\label{b3:prop:E-hecke-relations}
设 \(k\ge2\) 为偶整数，\(\Gamma=\mathrm{SL}_2(\mathbb Z)\)，\(T_n\) 为按\eqref{b3:eq:E-hecke-definition}定义的权 \(k\) Hecke 算子，其前因子为 \(n^{k-1}\)。在 \(M_k(\Gamma)\) 及其子空间 \(S_k(\Gamma)\) 上，对正整数 \(m,n\)，若 \((m,n)=1\)，则 \(T_mT_n=T_{mn}\)。对素数 \(p\) 与整数 \(r\ge1\)，
\[
 T_pT_{p^r}=T_{p^{r+1}}+p^{k-1}T_{p^{r-1}}.
\]
所有 \(T_n\) 彼此交换。
\end{proposition}
\begin{proof}
互素时，在两次系数公式中出现的公因子分别来自 \(m\) 与 \(n\)，彼此互素，乘积唯一对应 \(mn\) 的公因子；两次指标变化也相乘为 \(mn/d^2\)。对素数，写 \(a_x=0\) 当 \(x\) 不是非负整数，则
\((T_pa)_j=a_{pj}+p^{k-1}a_{j/p}\)。而
\((T_{p^r}a)_j=\sum_{h=0}^{\min(r,v_p(j))}p^{h(k-1)}a_{p^rj/p^{2h}}\)。
为核对截断处也没有遗漏，令 \(\nu=v_p(j)\)、\(C=p^{k-1}\)、
\(A=\min(r,\nu+1)\)、\(B=\min(r+1,\nu)\)。两次运算的第 \(j\) 个系数是
\[
\sum_{h=0}^{A}C^h a_{p^{r+1}j/p^{2h}}
+\sum_{h=1}^{B}C^h a_{p^{r+1}j/p^{2h}}.
\]
第二个和来自 \(C(T_{p^r}a)_{j/p}\)，把其旧指标加一即可；当 \(\nu=0\) 时它为空和。右侧算子 \(T_{p^{r+1}}+CT_{p^{r-1}}\) 给出同样两组项，只是两上界交换：第一组从0到 \(B\)，第二组从1到 \(A\)。零号项各出现一次，其他项的重数相同，故相等。对常数项，记 \(s_r=1+C+\cdots+C^r\)，则
\((1+C)s_r=s_{r+1}+Cs_{r-1}\)，这也验证了常数系数。故每个 \(T_{p^r}\) 为 \(T_p\) 的多项式，互素关系说明不同素数的算子交换，再分解 \(n\) 即得一般交换性。
\end{proof}

例如 \(T_2f\) 的展开不是把所有下标乘二，而是
\[
T_2f=\sigma_{k-1}(2)a_0+a_2q+
(a_4+2^{k-1}a_1)q^2+a_6q^3+
(a_8+2^{k-1}a_2)q^4+\cdots.
\]
偶数下标多出的第二项正是素数平方关系中的修正来源。取 \(r=1\) 得 \(T_4=T_2^2-2^{k-1}I\)，而不是 \(T_4=T_2^2\)。对 \(\Delta\)，因 \(T_2\Delta=-24\Delta\)，这个关系直接给 \(T_4\Delta=-1472\Delta\)。


\section{特征形式与 Euler 乘积}
\label{b3:sec:E03}

设 \(k\ge2\) 为偶整数，\(f\in S_k(\mathrm{SL}_2(\mathbb Z))\setminus\{0\}\)，并采用\eqref{b3:eq:E-hecke-definition}的归一化。若 \(f\) 对所有正整数 \(n\) 的 \(T_n\) 同时为特征向量，称为\term[hecke-tezheng-xingshi]{Hecke 特征形式}[Hecke eigenform]。若 \(T_nf=\lambda_nf\)，比较第一系数得 \(a_n=\lambda_na_1\)。所以 \(a_1=0\) 会迫使全部系数为零，与非零性矛盾。可以唯一把它正规化为 \(a_1=1\)，此时 \(\lambda_n=a_n\)。
\ProseParagraphBreak
本附录不以算子交换性直接断言一定存在完备的共同特征基；交换算子未必都可对角化。一般存在性还用 Petersson 内积下的自伴性等结构。当前已有一个明确例子：第二十三章证明 \(S_{12}\) 一维，由 \(\Delta=q-24q^2+252q^3-\cdots\) 张成，故每个 \(T_n\) 必在其中作用为标量，\(\Delta\) 是正规化特征形式。

\begin{theorem}[特征形式的乘积]\label{b3:thm:E-euler-product}
设 \(k\ge2\) 为偶整数，\(f\in S_k(\mathrm{SL}_2(\mathbb Z))\)，\(q=e^{2\pi\mathrm iz}\)，\(f(z)=\sum_{n\ge1}a_nq^n\)，\(a_1=1\)，且 \(f\) 为按\eqref{b3:eq:E-hecke-definition}归一的全部 \(T_n\) 的共同特征函数。则在 \(\operatorname{Re}s>k/2+1\) 内，级数与乘积绝对收敛，并有
\[
 L(s,f)=\sum_{n\ge1}a_nn^{-s}
       =\prod_p(1-a_pp^{-s}+p^{k-1-2s})^{-1}.
\]
\end{theorem}
\begin{proof}
Hecke 关系给 \(a_ma_n=a_{mn}\) 当 \((m,n)=1\)，以及
\(a_{p^{r+1}}=a_pa_{p^r}-p^{k-1}a_{p^{r-1}}\)。乘以 \(X^{r+1}\) 求和即得
\(\sum_{r\ge0}a_{p^r}X^r=(1-a_pX+p^{k-1}X^2)^{-1}\)。还需保证解析收敛。函数 \(y^{k/2}|f(z)|\) 对全模群不变，因为
\(\operatorname{Im}(\gamma z)=y/|cz+d|^2\)，恰好抵消 \(|f(\gamma z)|=|cz+d|^k|f(z)|\) 的因子。在基本域截去 \(y>Y\) 的部分后得到紧集，函数在那里有界；在尾部，\(f=qg(q)\)，其中 \(g\) 在足够小的闭圆盘有界，所以
\(y^{k/2}|f(z)|\le Cy^{k/2}e^{-2\pi y}\)，其界与 \(x\) 无关。每个点可移入基本域，故有统一的 \(|f(x+\mathrm{i}y)|\le Cy^{-k/2}\)。由
\[
 a_n=e^{2\pi ny}\int_0^1 f(x+\mathrm{i}y)e^{-2\pi\mathrm{i}nx}\,\mathrm{d}x
\]
取 \(y=1/n\)，得 \(|a_n|\le Cn^{k/2}\)。因此级数在 \(\operatorname{Re}s>k/2+1\) 绝对收敛，有限素数乘积可展开并按绝对收敛取极限，得到所示乘积。
\end{proof}
若把局部二次多项式分解为 \((1-\alpha_pX)(1-\beta_pX)\)，则 \(\alpha_p+\beta_p=a_p\)、\(\alpha_p\beta_p=p^{k-1}\)。本附录只证明这些代数关系，不假设 \(|\alpha_p|=|\beta_p|=p^{(k-1)/2}\)。后一个等模结论属于更深的理论，不能由乘积本身推出。

递推已经能给出许多可检查的局部计算。置 \(A=a_p\)、\(C=p^{k-1}\)，则
\[
a_{p^2}=A^2-C,\quad a_{p^3}=A^3-2CA,\quad
a_{p^4}=A^4-3CA^2+C^2.
\]
若两根不同，局部展开还给
\(a_{p^r}=(\alpha_p^{r+1}-\beta_p^{r+1})/(\alpha_p-\beta_p)\)；若重根为 \(\alpha_p\)，则为 \((r+1)\alpha_p^r\)。后一公式来自 \((1-\alpha_pX)^{-2}\)，所以不能把前一式的零分母直接代入。

\begin{proposition}[初始半平面的无零性]\label{b3:prop:E-initial-zero-free}
设 \(k\ge2\) 为偶整数，\(f\in S_k(\mathrm{SL}_2(\mathbb Z))\)，\(q=e^{2\pi\mathrm iz}\)，\(f(z)=\sum_{n\ge1}a_nq^n\)，\(a_1=1\)，且 \(f\) 为按\eqref{b3:eq:E-hecke-definition}归一的全部 \(T_n\) 的共同特征函数。则 \(L(s,f)=\sum_{n\ge1}a_nn^{-s}\) 在 \(\operatorname{Re}s>k/2+1\) 无零点。
\end{proposition}
\begin{proof}
固定该半平面内的 \(s\)。绝对收敛给
\(\sum_p\sum_{r\ge1}|a_{p^r}|p^{-r\operatorname{Re}s}<\infty\)。设局部和为 \(1+u_p\)，则 \(\sum_p|u_p|<\infty\)。局部递推与绝对收敛允许乘上二次多项式，给
\[
(1-a_pp^{-s}+p^{k-1-2s})(1+u_p)=1.
\]
所以各局部和都非零。除去有限个素数后 \(|u_p|<1/2\)，且 \(\sum_p\log(1+u_p)\) 绝对收敛，因此尾乘积是某个有限复数的指数，不为零；有限个头部因子也不为零。这就证明了整体无零性，而不仅是写出了一个形式乘积。
\end{proof}


\section{Mellin 变换与函数方程}
\label{b3:sec:E04}

对尖点形式，虚轴的两端都因变换律而指数衰减。这使 Mellin 积分比初始 Dirichlet 级数有更大的定义域。

\begin{theorem}[完成函数的方程]\label{b3:thm:E-functional-equation}
设 \(k\ge2\) 为偶整数，\(f\in S_k(\mathrm{SL}_2(\mathbb Z))\)，\(f(z)=\sum_{n\ge1}a_ne^{2\pi\mathrm inz}\)，最初在 \(\operatorname{Re}s>k/2+1\) 定义 \(L(s,f)=\sum_{n\ge1}a_nn^{-s}\)。不要求 \(f\) 是特征形式。则
\[
 \Lambda(s,f)=(2\pi)^{-s}\Gamma(s)L(s,f)
\]
延拓为整函数，且 \(\Lambda(s,f)=\mathrm{i}^k\Lambda(k-s,f)\)。
\end{theorem}
\begin{proof}
在绝对收敛区逐项积分得
\(\Lambda(s,f)=\int_0^\infty f(\mathrm{i}t)t^{s-1}\,\mathrm{d}t\)。当 \(t\to\infty\)，尖点条件给 \(f(\mathrm{i}t)=O(e^{-2\pi t})\)。矩阵 \(S\) 的变换律为 \(f(\mathrm{i}/t)=(\mathrm{i}t)^kf(\mathrm{i}t)\)，故当 \(t\downarrow0\)，\(f(\mathrm{i}t)=\mathrm{i}^{-k}t^{-k}O(e^{-2\pi/t})\)。积分及任意参数导数因此在全平面局部一致收敛，定义整函数。代换 \(t=1/u\) 后用变换律，得到
\[
 \Lambda(s,f)=\int_0^\infty f(\mathrm{i}/u)u^{-s-1}\,\mathrm{d}u
 =\mathrm{i}^k\int_0^\infty f(\mathrm{i}u)u^{k-s-1}\,\mathrm{d}u.
\]
这就是函数方程。再乘整函数 \(1/\Gamma(s)\) 得 \(L(s,f)\) 的整延拓。
\end{proof}
将 \(s\) 平移为 \(s+(k-1)/2\)，对称中心便从 \(k/2\) 移到 \(1/2\)。这解释了为何不同文献中的 \(L\)-函数虽然来自同一模形式，其 Gamma 因子和局部指数可能写得不同。

实际计算时，把积分在1处分割可得到不含小 \(t\) 端点的表达式：
\begin{equation}\label{b3:eq:E-split-mellin}
\Lambda(s,f)=\int_1^\infty f(\mathrm{i}t)
\bigl(t^{s-1}+\mathrm{i}^k t^{k-s-1}\bigr)\,\mathrm{d}t.
\end{equation}
第二项是对低端代换 \(t=1/u\) 后得到的。任意紧 \(s\) 集上，括号及各阶参数导数只增加固定幂和对数幂，均被 \(e^{-2\pi t}\) 控制；于是此式同时提供整延拓的直接计算式与逐次求导的依据。当 \(k\equiv0\pmod4\) 时，\(\Lambda(k/2+w,f)\) 是 \(w\) 的偶函数；当 \(k\equiv2\pmod4\) 时则为奇函数。相应地，前者中心的所有奇阶导数消失，后者所有偶阶导数消失，但函数方程本身并不强迫中心的一阶导数非零。

\ProseParagraphBreak
高层的类似计算需要明确一个附加对称。设 \(N\ge1\) 为整数、\(k\ge2\) 为偶整数，\(f\in S_k(\Gamma_0(N))\setminus\{0\}\)，并有 \(w\in\mathbb C\) 使 \(f|_kW_N=w f\)，其中 \(W_N=\begin{psmallmatrix}0&-1\\N&0\end{psmallmatrix}\)。因 \(W_N^2=-NI\)，偶权下 \(w^2=1\)。由定义
\(f(\mathrm{i}/(Nt))=wN^{k/2}(\mathrm{i}t)^kf(\mathrm{i}t)\)。为保留所有正规化因子，置 \(F_N(t)=f(\mathrm{i}t/\sqrt N)\)。上述变换律给
\[
F_N(1/t)=w\mathrm{i}^k t^k F_N(t).
\]
无穷远的尖点条件控制 \(t\to\infty\)，这一等式再控制 \(t\downarrow0\)。在初始收敛区逐项积分有
\[
\int_0^\infty F_N(t)t^{s-1}\,\mathrm{d}t
=\Gamma(s)\sum_{n\ge1}a_n\left(\frac{2\pi n}{\sqrt N}\right)^{-s}.
\]
然后代换 \(t=1/u\)，就得到
\[
 \Lambda_N(s,f)=\left(\frac{\sqrt N}{2\pi}\right)^s\Gamma(s)L(s,f),
 \qquad \Lambda_N(s,f)=w\mathrm{i}^k\Lambda_N(k-s,f).
\]
这是对已经满足 Fricke 特征关系的形式的条件性结论；没有该关系时，对称式会把 \(f\) 联系到另一尖点形式。

\section{自守表示与 Langlands 纲领概览}
\label{b3:sec:E05}

模形式的变换律可理解为某个离散群作用下的等变性。把函数转移到矩阵群上，再把平移生成的函数空间组织为表示，就进入\term[zishou-biaoshi]{自守表示}[automorphic representation]的语言。对 \(\mathrm{GL}_2\) 而言，经典模形式提供了许多这类表示的来源；无穷处记录权和解析类型，有限素数处记录局部 Hecke 特征值。
\ProseParagraphBreak
这种重组的实际意义在于把全局函数方程、Euler 因子与局部对称放在同一体系内。以本附录的例子看，一项 Fourier 展开的系数并非孤立数字：\(a_p\) 决定该素数的二阶递推，而所有素数的资料与 Mellin 积分必须相容。推广到更一般的群时，局部因子可能有更高次数，也可能需要不同的积分表示。
\ProseParagraphBreak
\term[langlands-gangling]{Langlands 纲领}[Langlands program]研究这些自守对象与 Galois 表示等算术对象之间的对应，以及不同群之间由对偶群映射所预期的传递。这里的“纲领”包含许多不同命题，不是一个已经在全部情形证明的总定理。椭圆曲线的模性、局部对应和高维的各种猜想需要各自的对象、假设与已知范围。本附录只建立理解这些问题所需的经典入口。进一步学习可参看 Bump《Automorphic Forms and Representations》\cite{Bump1997Automorphic}和 Iwaniec《Topics in Classical Automorphic Forms》\cite{Iwaniec1997Topics}。

\begin{chaptersummary}
Mellin 变换利用变换律产生延拓和函数方程；Hecke 特征关系利用有限求和产生乘法性和 Euler 乘积。这两种结构对同一组系数提出不同约束，其同时存在才形成模 \(L\)-函数的特殊性质。
\end{chaptersummary}
\BookExercises
\begin{exercise}\label{b3:ex:E-1}
若 \(f\in S_k\) 的第一系数为1且 \(T_pf=a_pf\)，证明 \(a_{p^2}=a_p^2-p^{k-1}\)。用于 \(\Delta\) 计算 \(a_4\)。
\end{exercise}
\begin{solution}
比较 \(T_pf\) 的第 \(p\) 个系数，左侧为 \(a_{p^2}+p^{k-1}a_1\)，右侧为 \(a_p^2\)。对 \(k=12,p=2,a_2=-24\)，得 \(a_4=576-2048=-1472\)。
\end{solution}
\begin{exercise}\label{b3:ex:E-2}
若 \(k\equiv2\pmod4\)，证明全模群尖点形式的 \(L(k/2,f)=0\)。
\end{exercise}
\begin{solution}
函数方程在中心给 \(\Lambda(k/2,f)=-\Lambda(k/2,f)\)，故为零。Gamma 因子在正数 \(k/2\) 有限非零，所以 \(L\) 同样为零。这里假定 \(k>0\)，如无非零该权尖点形式，结论自然仍成立。
\end{solution}
\begin{exercise}\label{b3:ex:E-3}
验证正规化的局部根 \(\alpha_p p^{-(k-1)/2}\)、\(\beta_p p^{-(k-1)/2}\) 的乘积为一，但仅凭此不能推出它们都在单位圆上。
\end{exercise}
\begin{solution}
乘积公式直接给一。数对 \((2,1/2)\) 已是反例，说明乘积为一允许模长互为倒数；还需要独立的根的估计。
\end{solution}
\begin{exercise}\label{b3:ex:E-4}
令 \(f(z)=\Delta(z)\)、\(g(z)=2^6\Delta(2z)\)，并取 \(W_2=\begin{psmallmatrix}0&-1\\2&0\end{psmallmatrix}\)。计算 \(f|_{12}W_2\)、\(g|_{12}W_2\)，并给出这两个函数张成空间中的两个 Fricke 特征向量。
\end{exercise}
\begin{solution}
定义与 \(\Delta(-1/z)=z^{12}\Delta(z)\) 给
\[
(f|_{12}W_2)(z)=2^6(2z)^{-12}\Delta(-1/(2z))
=2^6\Delta(2z)=g(z).
\]
同样，\((g|_{12}W_2)(z)=2^{12}(2z)^{-12}\Delta(-1/z)=f(z)\)。所以 \(f+g\)、\(f-g\) 分别具有特征值 \(+1\)、\(-1\)。它们都非零且线性独立，因为 \(f=q+O(q^2)\)、\(g=2^6q^2+O(q^4)\)。这给出了从交换的一对函数组成特征向量的具体计算。
\end{solution}
\begin{exercise}\label{b3:ex:E-5}
设 \(f\in S_k\)、\(k\equiv2\pmod4\)。把 \(\Lambda'(k/2,f)\) 写成从1到无穷远的积分，并证明 \(\Lambda''(k/2,f)=0\)。这些结论能否保证中心零点恰为单零点？
\end{exercise}
\begin{solution}
对式~\eqref{b3:eq:E-split-mellin} 求导，因 \(\mathrm{i}^k=-1\)，得
\[
\Lambda'(k/2,f)=2\int_1^\infty f(\mathrm{i}t)t^{k/2-1}\log t\,\mathrm{d}t.
\]
指数衰减保证求导合法。关于中心的奇性给二阶导数为零。上式积分仍可能为零，所以这些条件只限制零点阶数为奇数，并不能排除三阶或更高奇数阶；若 \(f=0\)，更不能谈有限零点阶数。
\end{solution}
\begin{exercise}\label{b3:ex:E-6}
利用 \(\Delta=q-24q^2+252q^3+\cdots\) 的 Hecke 特征关系，计算其第6、12个 Fourier 系数，不直接展开无穷乘积。
\end{exercise}
\begin{solution}
记 \(a_n\) 为 \(\Delta\) 的系数。互素关系给 \(a_6=a_2a_3=-24\cdot252=-6048\)。素数平方关系先给 \(a_4=(-24)^2-2^{11}=-1472\)，再因 \((3,4)=1\) 得 \(a_{12}=a_3a_4=-370944\)。若错误地使用完全乘法性，就会写出 \(a_{12}=a_3a_2^2\)，遗漏局部因子中的 \(2^{11}\) 项。
\end{solution}
